










\begin{derivation}[从\eqnrefs{eq:berezin-definition-dp68}到\eqnrefs{eq:grassmann-gaussian-determinant-dp68}]
Grassmann 高斯积分产生行列式的原因可以直接从“积分抽取最高次系数”看出。因为每个 $\theta_i$ 和 $\bar\theta_i$ 都平方为零，指数
\begin{equation*}
 \ee^{-\bar\theta M\theta}
 =\sum_{r=0}^{n}\frac{(-1)^r}{r!}(\bar\theta M\theta)^r
\end{equation*}
是有限多项式。测度含全部 $2n$ 个变量，所以积分后只有 $r=n$ 项可能留下。

展开 $(\bar\theta M\theta)^n$ 时，只要两个行指标或两个列指标重复，相应单项式就含某个 Grassmann 变量的平方而消失。因此非零项恰由 $1,\ldots,n$ 的排列标记。把变量重排回正文固定的测度次序会产生排列符号，而第一组指标的 $n!$ 种排列抵消展开中的 $1/n!$。于是
\begin{align*}
 \mathcal Z_G[0,0]
 &=\sum_{\sigma\in S_n}\operatorname{sgn}(\sigma)
 \prod_{i=1}^{n}M_{i\sigma(i)}\\
 &=\det M,
\end{align*}
最后一行正是行列式的 Leibniz 定义，从而得到正文\eqnrefs{eq:grassmann-gaussian-determinant-dp68}。
\end{derivation}

\begin{derivation}[从\eqnrefs{eq:grassmann-source-definition-dp68}到\eqnrefs{eq:grassmann-source-gaussian-dp68}]
加入源以后，Grassmann 高斯积分与普通高斯一样可以完成平方，只是变量和源的次序必须保持固定。令
\begin{equation*}
 \theta'=\theta-M^{-1}\xi,
 \qquad
 \bar\theta'=\bar\theta-\bar\xi M^{-1},
\end{equation*}
也就是 $\theta=\theta'+M^{-1}\xi$、$\bar\theta=\bar\theta'+\bar\xi M^{-1}$。代回指数并保持因子次序，有
\begin{align*}
 &-\bar\theta M\theta+\bar\xi\theta+\bar\theta\xi\\
 &=-\bar\theta'M\theta'
 -\bar\theta'\xi-\bar\xi\theta'-\bar\xi M^{-1}\xi
 +\bar\xi\theta'+\bar\xi M^{-1}\xi
 +\bar\theta'\xi+\bar\xi M^{-1}\xi\\
 &=-\bar\theta'M\theta'+\bar\xi M^{-1}\xi.
\end{align*}
Berezin 积分对这种常数 Grassmann 平移保持不变，因此
\begin{align*}
 \mathcal Z_G[\xi,\bar\xi]
 &=\ee^{\bar\xi M^{-1}\xi}
 \int\prod_i\dd\bar\theta_i'\dd\theta_i'\,
 \ee^{-\bar\theta'M\theta'}\\
 &=\mathcal Z_G[0,0]\ee^{\bar\xi M^{-1}\xi},
\end{align*}
即正文\eqnrefs{eq:grassmann-source-gaussian-dp68}。
\end{derivation}

\begin{derivation}[从\eqnrefs{eq:grassmann-source-gaussian-dp68}到\eqnrefs{eq:grassmann-inverse-kernel-dp68}]
正文已经固定了源的次序，因此二点收缩只需按同一约定求两次 Grassmann 导数。先对 $\bar\xi_i$ 取左导数，
\begin{equation*}
 \frac{\overrightarrow\partial\mathcal Z_G}{\partial\bar\xi_i}
 =\mathcal Z_G(M^{-1})_{ik}\xi_k.
\end{equation*}
再从右边对 $\xi_j$ 求导，并令源为零，指数因子回到 $1$，从而
\begin{equation*}
 \left.
 \frac{1}{\mathcal Z_G[0,0]}
 \frac{\overrightarrow\partial\mathcal Z_G}{\partial\bar\xi_i}
 \frac{\overleftarrow\partial}{\partial\xi_j}
 \right|_{0}
 =(M^{-1})_{ij}.
\end{equation*}
这就是正文\eqnrefs{eq:grassmann-inverse-kernel-dp68}。若交换源或导数的次序，会因 Grassmann 奇偶性出现额外负号，所以以后把这种带 Grassmann 源的高斯积分用于费米场时，必须同时固定源、导数和测度的次序。
\end{derivation}

\begin{derivation}[从\eqnrefs{eq:berezin-linear-transform-dp68}到\eqnrefs{eq:berezin-jacobian-dp21}]
普通变量的 Jacobian 来自体积元，而 Berezin 测度由最高次系数的归一化决定，所以变换方向恰好相反。先看一维 $\theta'=a\theta$。若写 $\dd\theta'=C\dd\theta$，则\eqnrefs{eq:berezin-definition-dp68} 要求
\begin{equation*}
 1=\int\dd\theta'\,\theta'
 =Ca\int\dd\theta\,\theta=Ca,
\end{equation*}
故 $C=a^{-1}$。

多维时由\eqnrefs{eq:berezin-linear-transform-dp68} 有
\begin{equation*}
 \theta_1'\cdots\theta_n'
 =A_{1i_1}\cdots A_{ni_n}\theta_{i_1}\cdots\theta_{i_n}.
\end{equation*}
只有 $(i_1,\ldots,i_n)$ 为一个排列时该项非零；重排到固定次序所产生的符号正好组成行列式，因此
\begin{equation*}
 \theta_1'\cdots\theta_n'=(\det A)\theta_1\cdots\theta_n.
\end{equation*}
若 $\dd^n\theta'=C\dd^n\theta$，最高次系数的 Berezin 归一化给出 $C\det A=1$，于是得到正文\eqnrefs{eq:berezin-jacobian-dp21}。
\end{derivation}
